\cotefile{123}
\scanpage{4}
Lissification du \unsure{morphisme|morphismes} de Kostant
% archive sheet number: 3
\scanpage{5}
$G$ schéma en groupes semi-simple \uline{adjoint} sur $S$.

$\unsure{\underline{g}|\underline{\mathfrak g}}$ Algèbre de Lie comme Module sur $\mathrm{Sch}_{/S}$, donc schéma vectoriel sur $S$

$\underline{g}^{rss}$ sous-schéma ouvert des éléments ``réguliers \inserted{semi-}simples'' i.e. de nullité minima (``réguliers'' dans SGA $3\ \overline{\underline{\mathrm{XIII}}}$)

$T$ ``le'' tore maximal d'un couple de Killing (tore sur $S$)

$\underline t$ son algèbre de Lie, regardé comme schéma vectoriel sur $S$

$W$ ``le'' groupe de Weyl, agissant sur $T$, donc sur $\underline t$.

$\unsure{J|I|\mathfrak J}=\underline t/W$ ; $\underline t$ est parfois noté $J'$, donc $J'\longrightarrow J$ revêtement (ramifié).

$\underline t^{rss}$ ouvert de $\underline t$ des éléments ``réguliers \inserted{semi-}simples''.

\inserted{Il est stable sous $W$.}

$J^{rss}=\underline t^{rss}/W$

$\unsure{\mathrm{Bor}|\underline{\mathrm{Bor}}}$ schéma des \deleted{drapeaux}\inserted{Borels}, \inserted{$\mathrm{Tor}$ schéma des tores maximaux,} $\mathrm{Kill}$ schéma des couples de Killing

$\mathscr B$ le ``Borel universel'' $\subset G_{\mathrm{Bor}}$. donc $\mathscr B/\mathscr B_u\simeq T_{\mathrm{Bor}}$

$\underline{\mathscr L}$ son algèbre de Lie, qui est un fibré vectoriel ($+$ autres structures) sur $\unsure{B|\mathrm{Bor}|\mathfrak B}$, $\underline{\mathscr L}/\underline{\mathscr L}_u\simeq\underline t_{\mathrm{Bor}}$. On a \unsure{hom|un morphisme} surjectif $\boxed{\unsure{q|p}\colon\underline{\mathscr L}\longrightarrow\underline g}$ et un morphisme $\boxed{\bar f'\colon\mathscr L\longrightarrow\underline t=J'}$

$\underline{\mathscr L}^{rss}$ sous-fibré des \unsure{élémts|éléments} ``rég. ss'' dans $\underline g_{\mathrm{Bor}}$,
$=q^{-1}(\underline g^{rss})$.

On a un isom can
\[
\underline{\mathscr L}^{rss}\simeq \mathrm{Kill}\times_S\underline t^{rss}.
\]
et un carré cartésien
\diagram{\begin{tikzcd}[ampersand replacement=\&,column sep=large]
\mathrm{Lie}(\mathcal T)^{rss}\arrow[r,"\sim"'] \& \underline g^{rss}\arrow[d] \& \underline{\mathscr L}^{rss}\simeq\mathrm{Kill}\times\underline t^{rs}\arrow[l]\arrow[d]\\
\& \mathrm{Tor} \& \mathrm{Kill}\arrow[l]
\end{tikzcd}}
\margin{où $\mathcal T$ est le tore maximal universel}

Il y a une opération naturelle de $W$ sur $\mathrm{Kill}$, faisant de $\mathrm{Kill}$ un torseur sur $W_{\mathrm{Tor}}$, d'où une opération de $W$ sur le produit fibré $\underline{\mathscr L}^{rss}$, qui sur $\mathrm{Kill}\times\underline t^{rs}$ se réduit à l'\unsure{act\textsuperscript{n}|action} \uline{diagonale}. Comme $\underline g^{rss}\simeq\mathscr L^{rss}/W$, on trouve un hom naturel
\[
\underline g^{rss}\longrightarrow J=\underline t^{rss}/W
\]
% archive sheet number: 4
\scanpage{6}
donnant naissance au diagramme
\diagram{\begin{tikzcd}[ampersand replacement=\&]
{\unsure{\underline g|\underline{\mathfrak g}}}^{rss}\arrow[d,"f^{rss}"'] \& \underline{\mathscr L}^{rss}\arrow[l,"{{q}^{rss}}"']\arrow[d,"{\widetilde f'{}^{rss}}"]\\
{\unsure{J|I|\mathfrak J}}^{rss}=\underline t^{rss}/W \& J'=\underline t^{rss}\arrow[l].
\end{tikzcd}}

Supposons que $W$ opère \deleted{\unsure{triv.|trivial.}}\uline{librement} sur $\underline t^{rss}$ ; \inserted{il revient ici au \unsure{m\textsuperscript{me}|même} qu'il opère \deleted{\unsure{fidèl.|fidèlement}} \unsure{libr\textsuperscript{t}|librement|fidèlement} \deleted{\unsure{sur les}} sur les fibres géom de $\underline t$,} ce qui est le cas au moins si les car. résiduelles sont toutes $\ne 2$. Alors le diagramme précédent est cartésien ; \deleted{\unsure{il a un sens}}

Je dis que le morphisme $\underline g^{\unsure{rs|rss}}\longrightarrow J^{rss}$ se prolonge en un morphisme $\underline g\longrightarrow J$ (évidemment de façon unique), de sorte qu'on aura le diagramme commutatif
\diagram{\begin{tikzcd}[ampersand replacement=\&,column sep=small]
\underline g\arrow[d,"f"'] \& \& \underline{\mathscr L}\arrow[ll,bend right=13,"q"']\arrow[dl,"{\widetilde f'}"]\\
J=\underline t^{rss}/W \& J'=\underline t^{\deleted{rss}}/\arrow[l] \&
\end{tikzcd}}
\margin{\uline{NB} L'anneau affine de $J$ sur $S$ est le sous-anneau des invariants de l'anneau affine de $\underline g$ sur $S$, et $f$ décrit \unsure{le|la} \unsure{pb|prop.} universel d'envoyer $\underline g$ de façon invariante dans un schéma relatif \uline{affine} (restriction \unsure{possiblement|probablement} inutile, du moins si on prend des schémas relatifs séparés)}

Dém. de l'existence de $f$. On peut se ramener au cas $S=\operatorname{Spec}\mathbb Z$, $G$ déployé. \deleted{Alors il suffit de prouver que $f^{rss},q^{rss}$ se prolongent, et} Comme $\underline g$ est normal, $J$ affine, il suffit de le voir sur la \unsure{fibres|fibre} générique. Là c'est connu (Kostant ....)

Introduisons
\[
\underline g'=\underline g\times_J J'
\]
on trouve donc un diagramme
\diagram{\begin{tikzcd}[ampersand replacement=\&,column sep=large,execute at end picture={\draw[thin,gray] (gprime) ellipse (.40cm and .53cm);\draw[thin,gray] (gprime.south) to[bend right=12] (ell.south east) to[bend right=67] (ell.north) to[bend right=9] (gprime.north);}]
\underline g\arrow[d,"f"'] \& |[alias=gprime]|\underline g'\arrow[l]\arrow[d,"f'"'] \& |[alias=ell]|\underline{\mathscr L}\arrow[l,"p"]\arrow[ll,bend right=16,"q"']\arrow[dl,"{\bar f'}"]\arrow[r] \& D\\
J \& J'\arrow[l] \& \&
\end{tikzcd}}
avec $p$ morphisme projectif, $\bar f'$ morphisme lisse ; dans le cas favorable où $W$ opère \deleted{\unsure{fidèl\textsuperscript{t}|fidèlement}}\inserted{librement} sur $J'^{rss}$, i.e. fidèlement sur les fibres géom. de $\underline t$, alors $p$ \deleted{induit} est un isom birationnel, plus précisément il induit un isom $\mathscr L^{rss}\simeq\underline g'^{rss}$, où $\underline g^{rss}=\underline g'|\underline g^{rss}$.
% archive sheet number: 5
\scanpage{7}
Par réduction au cas de la base $\mathbb Z$ ou $\mathbb Z[1/\unsure{n|2}]$, et par application du Main theorem, on trouve que le plus grand ouvert de ${\unsure{\underline g|\underline{\mathfrak g}}}'$ sur lequel $p$ est un isom est formé des $x'\in\underline g'$ tels que $p^{-1}(x)$ soit fini (ou même, contienne \inserted{un} élément isolé) i.e. l'image inverse de la partie \unsure{ouverte|ouvert} $\underline g^{\mathrm{reg}}\subset\underline g$, formée des points $x$ tels que $q^{-1}(x)$ soit fini (ou comme ...), i.e. des $x\in\underline g$ contenus dans un nb \uline{fini} seulement de sous-algèbres de Borel. En vertu de Springer PM n\textsuperscript{o}30, 5.3, \deleted{\unsure{il semble|il résulte}}\inserted{on voit} que cet ouvert \deleted{\unsure{conti|contient}} rencontre les parties unipotentes de tt\textsuperscript{e} fibre, et on \uline{prévoit} que l'on puisse en déduire qu'il rencontre tte fibre de $f\colon\underline g-\unsure{J|I|\mathfrak J}$.

\uline{Questions}

0) $W$ opère-t-il fidèl\textsuperscript{mt} sur $\underline t^{\unsure{rss|rs}}$ ? \uuline{Non} \uline{en car.} 2, \deleted{\unsure{les|cas}} $A_1,B_2$ ... peut-être d'autres ....

a) $J$ est-il régulier i.e. $J'$ plat sur $J$

b) $f$ est-il plat (impliquant a) ; en fait $f$ plat $\Longleftrightarrow$ $J$ régulier et $f$ équidimensionnel

\margin{OK pour les ``bons'' nombres premiers ?}
\margin{\deleted{\uline{NB} \unsure{la façon d'exprimer} \unsure{ceci} \unsure{l'isom} de l'anneau aff\textsuperscript{ine} de $\underline t/W$ sur \unsure{un} \uline{quotient} \unsure{de} \unsure{le sous-anneau} des invariants de l'anneau affine de $\underline g$ sur $S$, \unsure{$J$} \unsure{est quotient de} \unsure{$f(\underline g)$} ?}}

c) Les fibres de $f$ géom. \unsure{irréductibles|intègres} ? Sont elles chacune l'adhérence d'une orbite sous $G$, le complémentaire d'icelle dans la fibre étant de codim $\geq 2$ (ceci impliquant que les fibres de $f$ sont normales, \unsure{alors} $f$ donc $f'$ normal, donc $\underline g'$ normal \inserted{si base normale}, donc $\underline{\mathscr L}-\underline g'-\underline g$ factorisation de Stein)
\margin{\unsure{*)}}

d) Les points de lissité de $f$ sont ils les points réguliers (quasi-ponctuellement) i.e. ceux au dessus desquels $q$ est \deleted{\unsure{étale} (ou au moins quasi} fini ? Sont-ce aussi les points où le centralisateur \inserted{dans $G$} est de dimension minimum ?, soit $r$ ?

e) \deleted{Les orbites \inserted{unipotentes} de $G$ sur $\tilde G$, et $\underline g$ \unsure{sur}} Définir structure orbitale pour la Chevalley sur $\mathbb Z$, suivant une filtration naturelle de $\underline t$ ; cette structure orbitale devrait être homologue à la structure orbitale de $G$ opérant sur $\tilde G$ (le groupe simplement connexe) ; relations entre les centralisateurs correspondants à des orbites homologues ? / Question analogue pour $\mathscr B,\mathscr L$.

\margin{f) Reprendre théorie de Kostant : A-t-on section de l'ouvert des pts réguliers etc ---}

\margin{g) Étudier fibre de $p\colon\underline{\mathscr L}-\underline g'$, i.e. les composantes connexes de la variété des sous-alg. de Borel de $\underline g$ contenant un élément donné. Quelle est sa dimension (relation avec dim du \unsure{stabilisateur|fixeur|fixateur|l'orbite} )?}
% archive sheet number: 6
\scanpage{8}
Lissification du \unsure{morphisme|morphismes} de Steinberg
% archive sheet number: 7
\scanpage{9}
$G$ schéma en groupes réductif sur base $S$.

$D=\mathrm{Bor}(G)$, schéma des Borels.

$\mathscr B\longrightarrow D$ le ``Borel universel''\quad $\mathscr B\subset G_D=G\times_S D$ sous-schéma en \unsure{grpes|groupes}

$\mathscr B_u\subset\mathscr B$ la partie unipotente

$\mathscr B\to \unsure{\mathcal T|T}_D=\mathscr B/\mathscr B_u$ tore sur $D$.

On sait que $\mathcal T_D$ est muni d'une donnée de descente de $D$ à $S$, d'où un tore $T$ sur $S$ \deleted{\unsure{canonique}} avec isom
\[
\deleted{\unsure{\mathcal T}\simeq} T_D\simeq T\times_S D,\quad
\text{donc } T_D{}^\wedge=M_D,\quad\text{où }M=T^\wedge\quad\text{réseau sur }S.
\]

$\deleted{\unsure{\mathrm{Kt}|\mathrm{Kil}}}K=\mathrm{Kill}(G)$, schéma des couples de Killing \inserted{$(T,B)$}, $K=\underline{\mathrm{Tormax}}(\mathscr B)$

$K\longrightarrow D$ \deleted{\unsure{$f$|le}} morphisme structural

$\mathcal T_K\hookrightarrow\mathscr B_K$ le \deleted{tore maximal}\inserted{couple de Killing} universel \quad $T_K=T\times_S K$

$N_K=\operatorname{Norm}(\mathcal T_K,G_K)$

$\deleted{P_K=\unsure{\mathrm{Transl}|\mathrm{Trans}}}$

$W_K=N_K/T_K$ opère sur $T_K,M_K$

On sait que ces opérations proviennent d'opérations d'un groupe $W/S$ sur $T,M$ (avec sa structure radicielle)
\[
W_K\simeq W\times_S K
\]

\deleted{\illegible{1}}\inserted{Soit} $I=T/W$.

On définit un \inserted{unique} hom. (canonique), invariant par autom intérieurs (donc par \unsure{tt|tous} autom de $G$, quand on les fait opérer sur $I$ via autom extérieurs)
\[
G\xrightarrow{f} I
\]
\margin{Défini sans supposer $G$ simpl. connexe ? \uline{Oui} $f$ normal \inserted{ou au moins plat} \unsure{(en tous cas ?)|en tous cas ??} [Mais alors il faudrait que $I$ soit lisse/$S$ en tous cas !] \uuline{Non}}

La restriction de $f$ à $T_0$, si $T_0\subset B_0$, est donnée par
\[
f|T_0=\unsure{\text{hom.}|\text{mor.}|\text{morph.}}\ \text{can.}\quad T_0\longrightarrow T_0/W_0\simeq I.
\]
% archive sheet number: 8
\scanpage{10}
Soit $G'$ défini par \deleted{\unsure{le}}la condition d'avoir un diagramme cartésien
\diagram{\begin{tikzcd}[ampersand replacement=\&,column sep=large,execute at end picture={\draw (middle.north west) rectangle (rightnode.south east);}]
\unsure{\mathscr X|X|\mathfrak X}=G\arrow[d,"f"'] \& |[alias=middle]|\mathscr X'=G'\arrow[l]\arrow[d,"f'"]\arrow[dl,phantom,"\text{\uline{cart}}" description] \& |[alias=rightnode]|\mathscr B=\widetilde{\mathscr X}\arrow[l,"p"']\arrow[dl,"{\unsure{\widetilde f'|\widetilde f}}"]\\
I \& I'=T\arrow[l,"\pi"'] \&
\end{tikzcd}}
Utilisant le fait que $f$ est normal, donc $f'$ normal, on trouve que \deleted{\unsure{$f$}} $G'$ est normal.

On va définir un hom
\[
p\colon\mathscr B\longrightarrow G'
\]
i.e. des hom
\[
\widetilde f\colon\mathscr B\longrightarrow T,\quad
\deleted{\unsure{u}}\underline h\colon\mathscr B\longrightarrow G
\qquad\text{tels que }\pi\widetilde f=fh
\]

On prendra
\[
\left\{\begin{aligned}
\underline h(g,B)&=g\qquad \text{i.e. }\underline h\text{ est le composé }\mathscr B\subset G_D=G\times_S D\xrightarrow{\mathrm{pr}_1}G\\
\underline{\widetilde f}(g,B)&=\dot g\in B/B_u(-)\simeq T(-),\quad\text{i.e. }\widetilde f\text{ est le composé}\\
&\hspace{12mm}\mathscr B\longrightarrow\mathscr B/\mathscr B_u=T_D\simeq T\times_S D\xrightarrow{\mathrm{pr}_1}T
\end{aligned}\right.
\]

Pour voir que $\pi\widetilde f=fh$, il suffit de le voir sur l'ouvert dense \inserted{$\mathscr B^{rs}$} de $\mathscr B$ des \deleted{\unsure{po|pa}} $(g,B)$ avec $g$ régulier ss, \deleted{\unsure{où|si on}} \unsure{donc|d'où} $g\in T\subset B$, et alors c'est trivial.

\uline{NB}\quad $\left\{\begin{aligned}
&h\text{ est propre, donc }p\text{ est propre (on va voir qu'il est birationnel)}\\
&\widetilde f\text{ est lisse.}
\end{aligned}\right.$

$G^{\mathrm{reg}}$ l'ouvert des points réguliers de $G=$ ouvert des pts de lissité de $f$

$G'^{\mathrm{reg}}$ son image inverse dans $G'$

$\mathscr B^{\mathrm{reg}}$ son image inverse dans $\mathscr B$.
% archive sheet number: 9
\scanpage{11}
\uline{Théorème}\qquad (i) (Pour mémoire) $p\colon\mathscr B\longrightarrow G'$ est projectif,
\[
\widetilde f\colon\mathscr B\longrightarrow T\quad\text{est lisse.}
\]

(ii) $p$ induit
\[
\boxed{\mathscr B^{\mathrm{reg}}\overset{\mathrm{dfn}}{=}\mathscr B^{\deleted{\unsure{\mathrm{reg}}}}|G'^{\mathrm{reg}}
\xrightarrow{\sim}G'^{\mathrm{reg}}}
\]
et $G'^{\mathrm{reg}}$ est le plus grand ouvert au dessus duquel $p$ est un isom.

(iii) (Pour mémoire) $\mathscr B-\mathscr B^{\mathrm{reg}}$ est de codim $\geq 2$ \deleted{\unsure{$\mathscr B$}} dans $\mathscr B$ (fibre par fibres sur $S$ (ou sur $D$)

$G'-G'^{\mathrm{reg}}$ a toutes ses composantes irréductibles de codimension 3 (fibre par fibre sur $S$)

(iv) \inserted{(Si $S$ \deleted{\unsure{régulier}}\inserted{normal} \deleted{\illegible{1}} loc. noeth.)}
\diagram{\begin{tikzcd}[ampersand replacement=\&,column sep=small]
\mathrm{Pic}(\mathscr B)\arrow[r,"\sim"] \& \mathrm{Pic}(\mathscr B^{\mathrm{reg}}) \& \mathrm{Pic}(G'^{\mathrm{reg}})\arrow[l,"\sim"']\ \deleted{\unsure{\mathrm{Pic}(G')}}\\
\mathrm{Pic}(D)\arrow[u,"{\wr}"'] \& \&
\end{tikzcd}}
\inserted{[Si $T$ non tordu, \deleted{\unsure{ceci signifie simplement}} \unsure{resp. plus intrinsèquement|ou plus intrinsèquement} \unsure{généralement|génériquement} si pour tt point maximal $\unsure{\eta|s|\xi}$ de $S$, l'image $\Gamma$ de $\pi_1(S,\bar\eta)$ dans $\operatorname{Ext}(G_{\bar\eta})$ est telle que $H^1(\Gamma,M)=0$] on a une réciproque si $S$ régulier.}

et en faisceautisant sur $S$ (pour la top. \unsure{éte|ét.|étale}, disons)
\diagram{\begin{tikzcd}[ampersand replacement=\&,column sep=small]
\& \underline{\mathrm{Pic}}_{\mathscr B/S}\arrow[r,"\sim"] \& \underline{\mathrm{Pic}}_{\mathscr B^{\mathrm{reg}}/S} \& \underline{\mathrm{Pic}}_{G'^{\mathrm{reg}}/S}\arrow[l,"\sim"']\\
M\arrow[r]\&P\simeq\underline{\mathrm{Pic}}_{D/S}\arrow[u,"{\wr}"',"{\text{\deleted{\unsure{réciproque}}}}" near start,"{\text{\deleted{\unsure{le même groupe}}}}"' near end] \& \&
\end{tikzcd}}
\inserted{(injectif si $G$ semisimple)}

De plus, le cône ample pour $\underline{\mathrm{Pic}}_{\mathscr B/S}$ est le $=$ que pour $\underline{\mathrm{Pic}}_{D/S}$, qui est le cône des poids $P^+$. \deleted{\unsure{Et on a}}

\uline{Démonstration}\quad (i) et (iii) connus (pour (iii) ref : Steinberg ....)

(ii) On commence à voir, posant \deleted{$G^{rs}=\unsure{\text{ouv}}$}

$G^{rs}\subset G^{\mathrm{reg}}$ ouvert des pts de $G$ qui sont réguliers semi-simples
\[
\left.\begin{aligned}
G'^{rs}&\subset G^{\mathrm{reg}}\\
\mathscr B^{rs}&\subset\mathscr B^{\mathrm{reg}}
\end{aligned}\right|\quad\text{images inverses,}
\]
que
\[
\mathscr B^{rs}\xrightarrow{\sim} G'^{rs}
\]
D'autre part, on sait par Steinberg que les fibres de $\mathscr B^{\mathrm{reg}}\longrightarrow G'^{\mathrm{reg}}$ sont finies, d'ailleurs c'est un morphisme propre. Comme $G'^{\mathrm{reg}}$ est normal, on conclut par le Main Theorem.
\margin{On sait aussi par Steinberg que les fibres des autres pts sont de dim $>0$.}
% archive sheet number: 10
\scanpage{12}
(iv) Supposons \deleted{d'abord} $S$ \deleted{régulier}\inserted{normal}, \inserted{donc aussi tous les schémas écrits, \deleted{\unsure{sauf $G'$}}} \deleted{Alors $\mathrm{Pic}(\mathscr B)\xrightarrow{\sim}\mathrm{Pic}(\mathscr B^{\mathrm{reg}})$} $\mathrm{Pic}(\mathscr B)\longrightarrow\mathrm{Pic}(\mathscr B^{\mathrm{reg}})$ injectif résulte \inserted{\unsure{notions|questions|considérations} de profondeur et} du fait que $\mathscr B-\mathscr B^{\mathrm{reg}}$ est de codim $\geq 2$, \deleted{\unsure{et $\mathrm{Pic}(G'$}} \inserted{et $\mathrm{Pic}(D)\longrightarrow\mathrm{Pic}(\mathscr B^{\mathrm{reg}})$} résulte de considération fibre par fibre.

Le fait que $\mathrm{Pic}(D)\longrightarrow\mathrm{Pic}(\mathscr B)$ \deleted{soit un isom.}\inserted{\uline{surjectifs}} résulte de EGA $\underline{\mathrm{IV}}$ 21.4.13 (\unsure{Err|EGA} $\underline{\mathrm{IV}}$ 53), compte tenu du

\uline{Lemme}\quad $G$ groupe \deleted{\unsure{réductif (resp.)}}\inserted{semi-simple} simplement connexe sur corps $k$, $B$ un Borel, alors $\mathrm{Pic}(B)\deleted{\unsure{=\mathrm{Pic}(T)}}\inserted{=\mathrm{Pic}(B^{\mathrm{reg}})}=\mathrm{Pic}(T)=H^1(k,\underline M)$ \deleted{\unsure{si $k$}} \inserted{($T=B/B_u$, $\underline M=\widehat T$)}

Par loc. cit, utilisant \deleted{la relation} $B\longrightarrow T=B/B_u$, on est ramené au cas où $\mathrm{Pic}(T)=0$ (NB $B_u$ se dévisse en groupes additifs).

\margin{On a $=$ $\mathrm{Pic}(B)\simeq\mathrm{Pic}(T)$}

Or on sait que $T=\prod_{k'/k} T'$, $T'$ un tore sur une \deleted{\unsure{algèbre}}\inserted{extension} étale \inserted{galoisienne (gpe $\Gamma$)}, donc on a \deleted{\unsure{une équivalence de cat.|équiv. de cat.}}
\[
\deleted{\underline{\mathrm{Inv}}(T)\simeq\underline{\mathrm{Inv}}(T')\quad\text{donc}}
\]
\[
\deleted{\mathrm{Pic}(T)}\quad
0\longrightarrow H^1\bigl(\Gamma,\underbrace{H^0(T',\underline O_{T'}^{\times})}_{M\times k'^{*}}\bigr)
\longrightarrow\mathrm{Pic}(T)\longrightarrow H^0\bigl(\Gamma,\underbrace{H^1(T',\underline O_{T'}^{\times})}_{=0}\bigr)
\]
\deleted{\unsure{partie adjointe est nulle} \unsure{au sens entier}}

donc on trouve
\[
\mathrm{Pic}(T)=H^1(\Gamma,M)\times\underset{\substack{=0\\\mathrm{H.}\,90}}{H^1(\Gamma,k'^{\times})}=H^1(\Gamma,M)=H^1(k,\underline M)
\]
Si $P$ est tordu, cela peut être $\ne 0$ ; si $B$ est la fibre générique de $\mathscr B$ sur $D_{\unsure{s|S|{}}}$ ($S$ irréd), alors cela implique donc que l'hom $\mathrm{Pic}(D)\longrightarrow\mathrm{Pic}(\mathscr B)$ n'est pas surjectif. L'injectivité est claire, car $\mathscr B$ a une section sur $D$ (la section nulle). En faisceautisant pour la topologie étale, on trouve donc le diagramme d'isoms sur les $\underline{\mathrm{Pic}}$. / Comme $\mathscr B\longrightarrow D$ est affine, et a une section, il s'ensuit qu'un faisceau inv. sur $D$ est ample ssi son image inverse sur $\mathscr B$ l'est.
\margin{indép. de l'hypothèse sur $S$ !} % slanted

\deleted{Que peut-on faire dans le cas non régulier ? $\Rightarrow\underline{\mathrm{Pic}}$ \unsure{$\subset$} a) L'argument établissant $\mathrm{Pic}(D)\xrightarrow{\sim}\mathrm{Pic}(\mathscr B)$ \unsure{marche|et} dès que $S$ est \uline{normal}. De $=$ pour $\mathrm{Pic}(D)\longrightarrow\mathrm{Pic}(\mathscr B^{\mathrm{reg}})$. L'injectivité t\textsuperscript{js}. \unsure{directement|décentrement|décentement|évidemment} !}
% archive sheet number: 11
\scanpage{13}
\uline{Remarques} a) Si $T$ peut être tordu, on ne peut plus affirmer que l'hom injectif
\[
\mathrm{Pic}(D)\longrightarrow\mathrm{Pic}(\mathscr B)
\]
soit surjectif. Mais (si $S$ \deleted{\unsure{est noeth.})} n'a qu'un nb fini de composantes connexes) on voit que le \uline{conoyau est fini} en fait on a \deleted{pour} $S$ étant de \unsure{dbl|ctible} une suite exacte
\[
\begin{aligned}
0-\mathrm{Pic}(D)-\mathrm{Pic}(\mathscr B)&\longrightarrow H^1(\Gamma,M_{\bar{\unsure{\xi|\zeta}}}),\\[-1pt]
&\hspace{8mm}\Vert\\[-2pt]
&\mathrm{Pic}(\mathscr B_\eta)\xrightarrow{\sim}\mathrm{Pic}(T_\eta)\xleftarrow{\sim}\mathrm{Pic}(T_\xi)
\end{aligned}
\]
où $\eta$ pt gén. de $D$\quad $\xi$ pt gén. de $S$

\deleted{\illegible{1}} où $\Gamma$ est l'image de $\pi_1(S,\bar\xi)$ dans $\operatorname{Aut}(M_{\bar\xi})$.

Si $S$ est régulier, \deleted{\unsure{et alors}} on peut mettre un zéro à droite, i.e. la condition $H^1(\Gamma,M_{\bar\xi})=0$ est néc. et suff. pour que l'on ait bijectivité $\mathrm{Pic}(D)-\mathrm{Pic}(\mathscr B)$ \deleted{\unsure{et}} \deleted{\unsure{elle implique en tous cas inject.}}\inserted{ou même} bijectivité $\mathrm{Pic}(D)\simeq\mathrm{Pic}(\mathscr B^{\mathrm{reg}})$ [car $\mathrm{Pic}(\mathscr B_\eta)\longrightarrow\mathrm{Pic}(\mathscr B_{\unsure{p|\eta|\xi}}^{\mathrm{reg}})$ est \inserted{en tous cas} \unsure{bijectif|injectif|surjectif} pour des raisons de codim $\geq 2$ et de régularité]

b) L'injectivité de $\mathrm{Pic}(D)-\mathrm{Pic}(\mathscr B)$ et la caractérisation des éléments de $\mathrm{Pic}(D)$ amples/$S$ comme ceux dont l'image inverse sur $\mathscr B$ sont amples /$S$, restent vrais sans hyp. de normalité sur $S$.

c) Supposons $S$ normal \inserted{et $H^1(\xi,M)=0$ p.ex. $T$ non tordu}. Alors
\[
\mathrm{Pic}(S)\xrightarrow{\sim}\mathrm{Pic}(G')
\]

En effet, on \deleted{sait que $\mathrm{Pic}(G')$} a injectivité car $G'$ a une section sur $S$ (section unité). Reste à prouver que si $\unsure{\mathscr F|\varphi|\xi}\in\mathrm{Pic}(G')$ induit 0 sur la section unité, alors $\mathscr F$ est nul. Or
\[
\mathrm{Pic}(G')\hookrightarrow\mathrm{Pic}(G'^{\mathrm{reg}})\bigl[\simeq\mathrm{Pic}(\mathscr B^{\unsure{\mathrm{reg}|\prime\mathrm{reg}}})\xleftarrow{\sim}\mathrm{Pic}(\mathscr B)\bigr]
\]
inj. pour raisons de profondeur, donc il suffit \deleted{que} de voir que l'image \deleted{\unsure{inv|inverse}} de $\mathscr F$ dans $\mathrm{Pic}(\mathscr B)$ est nulle. Mais il suffit de voir que la restriction de cette dernière à la section unité de $\mathscr B$ sur $D$ est nulle. Mais $D\longrightarrow G'$
% archive sheet number: 12
\scanpage{14}
se factorise par la section unité de $G'$ sur $S$ ! En fait, sans hyp. sur $T$, on obtient toujours une suite exacte
\diagram{\begin{tikzcd}[ampersand replacement=\&,row sep=small]
G'^{\mathrm{reg}}\arrow[d,hook] \& \mathscr B^{\mathrm{reg}}\arrow[l]\arrow[d,hook]\\
G'\arrow[dd] \& \mathscr B\arrow[l]\\
\& D\arrow[u,hook]\arrow[ul]\arrow[dl]\\
S\arrow[uu,bend right=17] \&
\end{tikzcd}}
\[
0-\mathrm{Pic}(S)\longrightarrow\mathrm{Pic}(G')\longrightarrow
\underset{\substack{\Vert\\\text{groupe fini}}}{H^1(\xi,M)}
\]

d) On a \deleted{$\mathrm{Pic}(G)$} ($S$ normal \unsure{\& irréd. de pt gén.|irréd. de point générique} $\xi$)
\[
0-\mathrm{Pic}(S)\underset{\substack{\text{injectif à cause}\\\exists\ \text{section unité}}}{\longrightarrow}\mathrm{Pic}(G)\longrightarrow\mathrm{Pic}(G_\xi)
\]
groupe \deleted{\unsure{fini t\textsuperscript{jours}|fini tors.}} qui est nul si $G_\xi$ est simplement connexe et déployé.

\deleted{\unsure{e)}} D'autre part \deleted{si $S$ régulier} on a $\mathrm{Pic}(G)\xrightarrow{\mathrm{surj}}\mathrm{Pic}(G^{\mathrm{reg}})$ (\inserted{\unsure{sous|sans} hyp. de normalité,} car on a codim $\geq3$ fibre par fibre \inserted{sur $S$} de $G-G^{\mathrm{reg}}$), donc on trouve
\[
\operatorname{Coker}\bigl(\mathrm{Pic}(S)-\mathrm{Pic}(G^{\mathrm{reg}})\bigr)
\quad\text{est fini (et $=$ nul si $G_\xi$ est simplement connexe et déployé)}
\]

Cela montre \deleted{bien} qu'on ne peut trouver une résolution $\widetilde X$ de $f$ analogue à celle de $f'$, sans toucher à $G^{\mathrm{reg}}$ (l'image inverse de $G-G^{\mathrm{reg}}$ dans $\widetilde X$, y étant de codim $\geq2$ fibre par fibre ...)

\deleted{e)}
% archive sheet number: 13
